% ============================================================
% Baseball2Vector — CVPR Workshop (CVsports) style draft
% ------------------------------------------------------------
% HOW TO USE ON OVERLEAF:
%  1) Start from the official CVPR template project
%     (search "CVPR 2026" in Overleaf templates), which provides
%     cvpr.sty and the review/final options.
%  2) Replace the main .tex with this file.
%  3) Copy the three figures from outputs/figures/ into a
%     figures/ folder in the Overleaf project:
%       - fig1_zscore_area_vs_war.png
%       - fig2_radar_witt_vs_rutschman.png
%       - roi_coefficients.png
%  All numbers in this draft are verified against
%  data/processed/v3_results.csv and outputs/tables/ in the repo.
% ============================================================

\documentclass[10pt,twocolumn,letterpaper]{article}

% --- CVPR template package (from the official template project) ---
\usepackage[pagenumbers]{cvpr}   % use [review] for the review version

\usepackage{graphicx}
\usepackage{amsmath,amssymb}
\usepackage{booktabs}
\usepackage{multirow}
\usepackage{xurl}
\usepackage[pagebackref,breaklinks,colorlinks]{hyperref}

\def\paperID{****}
\def\confName{CVPR Workshop (CVsports)}
\def\confYear{2026}

\begin{document}

%%%%%%%%% TITLE
\title{Baseball2Vector: An Interpretable Five-Tool Representation of MLB Position Players}

\author{Jaeseok Choi\\
I2SLAB, Sungkyunkwan University\\
{\tt\small jaeseok.choi@skku.edu}
}

\maketitle

%%%%%%%%% ABSTRACT
\begin{abstract}
Modern baseball evaluates players through a proliferation of statistics---raw physical measurements, sabermetric composites such as WAR and wRC+, and the traditional 20--80 scouting scale---yet each captures only a partial view, and scalar composites compress qualitatively different profiles onto a single axis. We propose \emph{Baseball2Vector}, a framework that represents each MLB position-player season as a five-dimensional vector across Power, Contact, Plate Discipline, Defense, and Speed. These dimensions are inspired by, but are not identical to, the traditional five scouting tools; their values are reported on a bounded 20--80-formatted relative scale and should not be interpreted as professional scout grades. Using FanGraphs statistics for 2{,}309 hitter-seasons (2021--2025, $\geq$100 PA), we compare per-dimension Z-score averaging, per-dimension PCA, and an exploratory joint variational autoencoder, and analyze the area of the resulting radar-chart pentagon as a geometric summary. Z-score averaging shows the strongest descriptive alignment with established value metrics (Pearson $r$ of 0.81, 0.65, and 0.63 against WAR, wRC+, and OPS, respectively). Because the inputs overlap with components of those metrics, these correlations are internal-consistency checks rather than independent validation. An observational regression of concurrent year-over-year dimension and value changes attains a reported five-fold cross-validated $R^2$ of 0.71; we interpret its synthetic $+5$-grade perturbations as model-sensitivity estimates, not causal counterfactuals or development prescriptions. Baseball2Vector is therefore positioned as an interpretable interface for examining player profiles, rather than as a replacement for value metrics or scouting evaluation.
\end{abstract}

%%%%%%%%% BODY
\section{Introduction}
\label{sec:intro}

Evaluating and comparing player ability is central to modern baseball. For front offices, accurately gauging a player's ability underpins the decisions that matter most: which prospects to draft or sign, and which players to deploy in which roles. For players themselves, an objective account of where their strengths and weaknesses lie is a prerequisite for targeted development. And for fans, much of the enjoyment of the sport comes from understanding \emph{how} a player is good---which tools carry them---and from counterfactual imagination: ``what if this player added power?'' In all three cases, what is needed is not merely a ranking but a \emph{representation}: a compact, interpretable description of a player that supports comparison, communication, and reasoning.

Existing answers are only partial. Sabermetric composites such as WAR and wRC+ condense a player's contribution into a single number, which is invaluable for ranking but discards profile information: a slugger with poor contact skills and a well-rounded defender with modest power can receive similar scores while being entirely different players. Conversely, raw tracking statistics preserve detail but can overwhelm interpretation. The traditional 20--80 scouting scale sits between these extremes and provides a widely shared vocabulary in the sport, but it is typically assigned through expert observation and projection rather than derived reproducibly from public season statistics \cite{mcdaniel2014}.

We propose \emph{Baseball2Vector}, which represents each position-player season as a five-dimensional vector across Power, Contact, Plate Discipline, Defense, and Speed, computed from public statistics and displayed on a bounded 20--80-formatted relative scale. This is an adapted skill taxonomy rather than a reconstruction of the canonical five scouting tools. The vector can be visualized as a pentagon; its area is analyzed as an exploratory geometric summary, while the vector itself retains the dimension-level profile that the area discards.

Our contributions are as follows:
\begin{itemize}
  \item We formalize a data-driven mapping from public season statistics to a five-dimensional, 20--80-formatted relative profile, and compare per-dimension Z-score averaging, per-dimension PCA, and an exploratory joint variational autoencoder.
  \item We characterize the radar-chart pentagon area as an ordering-sensitive geometric summary and report its descriptive alignment with WAR, wRC+, and OPS over 2{,}309 hitter-seasons from 2021--2025.
  \item We show that the simplest encoding (Z-score averaging) has the strongest internal alignment with the selected value metrics ($r=0.81$ vs.\ WAR), while separating this shared-input sanity check from external construct validation.
  \item We use year-over-year dimension-delta regression to quantify observational associations and local model sensitivity. Synthetic perturbations are treated as hypothesis-generating diagnostics rather than causal estimates of player development.
\end{itemize}

\section{Related Work}
\label{sec:related}

\paragraph{Measurement-level statistics.}
Baseball Savant\footnote{\url{https://baseballsavant.mlb.com}} publicly releases a large volume of in-game physical measurements (exit velocity, sprint speed, arm strength, and so on) captured by the Statcast system. These measurements are enormously useful to teams, players, and fans, but they are \emph{measured outcomes}: interpreting them requires at least one further stage of aggregation and normalization before they speak to player ability.

\paragraph{Sabermetric value composites.}
FanGraphs\footnote{\url{https://www.fangraphs.com}} and Baseball Reference\footnote{\url{https://www.baseball-reference.com}} publish composite value metrics such as WAR and wRC+ that summarize a player's total contribution. These have become the lingua franca of player valuation, but they carry known limitations: different publishers use different formulas, so the ``same'' metric takes different values across sources; and by construction a single scalar cannot distinguish players who reach the same value through different skill profiles.

\paragraph{The 20--80 scouting scale.}
Long before Statcast, scouts graded players on a 20--80 scale centered at 50, with each 10-point step conventionally corresponding to approximately one standard deviation. The canonical position-player tools are Hit, Power, Run, Field, and Throw; plate discipline is commonly treated as a component of the Hit tool rather than a separate sixth tool \cite{mcdaniel2014}. The scale preserves profile information that value composites discard, but grades are traditionally assigned through expert observation and projection. Baseball2Vector borrows the visual range, not the full semantics or validity of professional scouting grades.

\paragraph{Baseball player representations.}
Prior work has learned latent representations of baseball players from plate-appearance interactions. Alcorn's \emph{batter2vec/pitcher2vec} predicts at-bat outcomes from neural player embeddings \cite{alcorn2016}, while Heaton and Mitra learn time-varying 72-dimensional player-form representations from sequences of Statcast events \cite{heaton2021}. These approaches emphasize predictive or contextual embeddings. Our narrower goal is an explicitly named, season-level profile whose dimensions can be inspected directly; accordingly, interpretability is prioritized over embedding capacity.

\paragraph{Vector representations in NLP.}
In natural language processing, Word2Vec demonstrated that words can be embedded in a semantic vector space in which geometric proximity reflects semantic similarity and, famously, arithmetic on vectors captures analogical relations (\emph{king} $-$ \emph{man} $+$ \emph{woman} $\approx$ \emph{queen}) \cite{mikolov2013}. This result established distributed vector representations as a general tool for turning discrete entities into computable objects.

\paragraph{This work.}
Inspired by these threads, Baseball2Vector uses public season-level statistics to encode observed player performance as a low-dimensional, interpretable vector. Unlike black-box player embeddings, each dimension is assigned a baseball-relevant name; unlike scalar composites, the representation preserves profile structure. Synthetic modifications of a vector are used only to probe the fitted model's sensitivity, not to claim that an intervention on one skill would cause the predicted change in value.

\section{Methodology}
\label{sec:method}

\subsection{Tool Selection}
The canonical five tools for position players are Hit, Power, Run, Field, and Throw \cite{mcdaniel2014}. Public FanGraphs season statistics do not provide a satisfactory, consistently available measure of throwing-arm ability for hitters across the entire study window. We therefore define an \emph{adapted five-dimensional taxonomy}: \textbf{Power}, \textbf{Contact}, \textbf{Plate Discipline}, \textbf{Defense}, and \textbf{Speed}. This taxonomy omits the traditional Throw tool and separates plate discipline from the broader Hit tool. We consequently refer to these quantities as \emph{dimensions} rather than canonical scouting tools, and do not interpret them as direct substitutes for professional present or future grades.

\subsection{Data Collection and Preprocessing}
\label{sec:data}
We collect season-level statistics from FanGraphs (via the \texttt{pybaseball} library\footnote{\url{https://github.com/jldbc/pybaseball}}) for all MLB hitters with at least 100 plate appearances in each of the 2021--2025 seasons, yielding 2{,}309 player-season rows. Each row is keyed as \emph{Player (Season)} so that repeated seasons from the same player remain identifiable rather than being treated as independent players. A data audit found one player-name collision affecting this key---two distinct major leaguers both named Max Muncy each qualified in 2025---which we resolved by relabeling the two rows using performance context (PA, WAR, wRC+) rather than dropping either qualifying season; no other name collisions, duplicate rows, or out-of-range statistics were found in the 2{,}309-row dataset. Constituent statistics are grouped by dimension as follows:
\begin{itemize}
  \item \textbf{Contact}: Contact\%, K\%, AVG, xBA, SwStr\%
  \item \textbf{Power}: ISO, SLG, HardHit\%, Barrel\%, maxEV, EV, HR/FB
  \item \textbf{Speed}: Spd, BsR, UBR, wSB
  \item \textbf{Plate Discipline}: BB\%, O-Swing\%, Swing\%, BB/K
  \item \textbf{Defense}: Def, Fld
\end{itemize}
Statistics for which lower is better (K\%, SwStr\%, O-Swing\%, Swing\%) are sign-flipped before aggregation so that higher always means better.

Because rate statistics computed over few plate appearances are noisy, we apply empirical-Bayes-style shrinkage before scoring: each rate statistic is pulled toward its season league average using a pre-specified prior strength expressed in PA equivalents (e.g., 60 PA for K\%, 120 PA for BB\%, and 500 PA for AVG). These are heuristic hyperparameters motivated by published stabilization-point summaries \cite{fangraphsSampleSize}, rather than estimated causal or population priors. A post hoc sensitivity check (Sec.~\ref{sec:results}) finds the headline alignment correlations move by at most 0.003 under substantially stronger shrinkage than currently applied, suggesting the specific thresholds are not a load-bearing choice, at least in that direction. Statistics that are already composite measures (e.g., Spd and Def) are left untouched. All standardization is then computed \emph{within season}---each statistic is Z-scored against that season's population---to respect era and run-environment differences. For predictive evaluation, these preprocessing parameters must be fitted using training data only; the present pooled alignment analysis is descriptive and does not constitute an out-of-sample test.

\subsection{Encoding Schemes}
Each dimension $t$ has a set of constituent statistics $S_t$. We compare three ways of collapsing $S_t$ into a single dimension score.

\paragraph{Z-score averaging.}
Each statistic in $S_t$ is standardized to zero mean and unit variance over the per-season player population; the dimension score is the mean of these Z-scores. Because several inputs within and across dimensions are correlated or algebraically related (e.g., AVG and SLG, or K\% and BB/K), the resulting averages are transparent composites rather than estimates of independent latent abilities.

\paragraph{PCA.}
The statistics in each $S_t$ co-vary as facets of a related performance construct (e.g., ISO, SLG, and Barrel\% all reflect power). We apply PCA within each dimension group and take the first principal component as the dimension score, capturing its dominant shared axis of variation.

\paragraph{JointVAE.}
Per-dimension encoders cannot capture cross-dimension structure. We therefore also train a single variational autoencoder over all constituent statistics jointly \cite{kingma2014}: a shared MLP encoder maps the full statistic set to a 5-dimensional latent with a full posterior covariance parameterized through its Cholesky factor ($\Sigma = LL^{\top}$), and five group-specific decoders reconstruct the input groups. Latent dimensions are assigned post hoc to named dimensions using correlations with anchor statistics (e.g., ISO for Power), with signs flipped where needed. This procedure does not guarantee disentanglement or unique semantic identification; accordingly, the JointVAE is treated as an exploratory representation baseline. Posterior off-diagonal covariance describes dependence in the approximate latent posterior and is not interpreted directly as a causal or population-level tool trade-off. Results are reported for a fixed seed; stability across random seeds remains to be assessed.

\paragraph{20--80 scaling.}
Whatever the encoder, dimension scores are mapped to a bounded 20--80-formatted range by an affine min--max transformation over the pooled 2021--2025 population, so that the observed extremes anchor 20 and 80. The result is a vector $\mathbf{v} = (v_{\text{Con}}, v_{\text{Pow}}, v_{\text{Spd}}, v_{\text{Def}}, v_{\text{Dis}}) \in [20,80]^5$. This transformation is not the classical scouting convention of mean 50 and approximately 10 points per standard deviation. It is cohort-dependent: adding a new extreme observation can change every player's grade. We therefore describe the outputs as \emph{20--80-formatted relative scores}, not scouting grades.

\subsection{Pentagon Area as a Geometric Summary}
For visualization, $\mathbf{v}$ is drawn as a radar chart: five axes at equal $72^\circ$ intervals. We analyze the area of the resulting pentagon as a compact geometric descriptor, computed by decomposing it into five triangles:
\begin{equation}
A(\mathbf{v}) = \tfrac{1}{2}\sin\!\left(\tfrac{2\pi}{5}\right) \sum_{i=1}^{5} v_i\, v_{i+1},
\label{eq:area}
\end{equation}
with indices cyclic. Because Eq.~\ref{eq:area} is a sum over adjacent pairs, the area depends on the ordering of axes. We quantify this sensitivity in Sec.~\ref{sec:results}. Area is therefore not proposed as an invariant overall-value metric; it is a lossy visualization-derived summary whose behavior helps motivate retaining the full vector.

\subsection{Alignment Protocol}
As an internal-consistency check, we compute Pearson correlations between pentagon area and three established metrics---WAR, wRC+, and OPS---over all 2{,}309 season-rows, separately for each encoding scheme. These are pooled, descriptive correlations. Several constituent statistics also contribute directly or indirectly to the comparison metrics; moreover, repeated seasons from the same player are not independent observations. The analysis therefore does not establish external construct validity or predictive performance.

\subsection{Dimension-Delta Regression and Model Sensitivity}
\label{sec:roi-method}
To study how concurrent changes in the representation relate to changes in established metrics, we construct a transition dataset of all consecutive-season pairs of the same player (1{,}427 pairs), with features $\Delta\mathbf{v} \in \mathbb{R}^5$ and targets $\Delta$wRC+, $\Delta$WAR, and $\Delta$OPS. We fit (i) a linear regression on the five deltas and (ii) a random forest on ten features (the five starting levels plus the five deltas). Both are evaluated under three cross-validation schemes: (i) pooled five-fold cross-validation over transition rows, as overlapping transitions from the same player may otherwise share player-specific information across folds; (ii) player-grouped five-fold cross-validation (\texttt{GroupKFold} keyed on player identity, so no player's transitions appear in both a fold's training and test split), which additionally refits the 20--80 scaling of Sec.~\ref{sec:method} on each training fold only; and (iii) a forward-time holdout that trains on all transitions ending in 2023 or earlier and evaluates only on 2024$\to$2025 transitions, the closest available approximation to a genuinely prospective forecast. We probe local model sensitivity by applying a synthetic $+5$-grade perturbation to each dimension in turn and recording the random forest's change in predicted $\Delta$wRC+. This perturbation is not a causal intervention, and ``return on investment'' is not claimed because neither intervention cost nor causal identification is modeled.

\section{Results}
\label{sec:results}

\subsection{Quantitative Alignment}
Table~\ref{tab:corr} reports Pearson correlations between pentagon area and established value metrics for the three encodings, and Fig.~\ref{fig:scatter} shows the corresponding scatter for the Z-score encoding against WAR.

\begin{table}[t]
\centering
\caption{Pearson correlation between pentagon area and established value metrics, by encoding scheme (2021--2025, $n=2{,}309$ season-rows). Brackets are 95\% confidence intervals from a player-level cluster bootstrap (2{,}000 resamples; players, not season-rows, are the resampling unit, since a player's several seasons are not independent draws). A paired bootstrap on the same resamples confirms Z-score's advantage over PCA, though only 0.01--0.02 in magnitude, is statistically significant (95\% CI on the difference excludes 0) for all three metrics; the larger gap to JointVAE (0.05--0.16) is significant a fortiori.}
\label{tab:corr}
\small
\begin{tabular}{lccc}
\toprule
Encoding & WAR & wRC+ & OPS \\
\midrule
Z-score averaging & \textbf{0.81} {\scriptsize[0.78, 0.83]} & \textbf{0.65} {\scriptsize[0.61, 0.68]} & \textbf{0.63} {\scriptsize[0.59, 0.66]} \\
PCA               & 0.79 {\scriptsize[0.77, 0.82]} & 0.63 {\scriptsize[0.60, 0.67]} & 0.61 {\scriptsize[0.57, 0.65]} \\
JointVAE          & 0.74 {\scriptsize[0.71, 0.77]} & 0.51 {\scriptsize[0.47, 0.55]} & 0.49 {\scriptsize[0.44, 0.53]} \\
\bottomrule
\end{tabular}
\end{table}

\begin{figure*}[t]
\centering
\includegraphics[width=0.75\linewidth]{figures/fig1_zscore_area_vs_war.png}
\caption{Pentagon area (Z-score encoding) vs.\ WAR over all 2{,}309 season-rows ($r=0.81$). The pooled correlation is descriptive and partly reflects statistics shared by the representation and WAR.}
\label{fig:scatter}
\end{figure*}

Four observations follow. First, Z-score averaging has the strongest descriptive alignment with all three metrics, with PCA close behind and the JointVAE lower. This result favors the transparent baseline for the present purpose, but it does not show that learned encodings are generally inferior. Second, although Z-score's and PCA's marginal confidence intervals in Table~\ref{tab:corr} overlap substantially, a paired bootstrap on the same player resamples---which accounts for the fact that all three encodings are computed on the same players and are themselves highly correlated with each other---shows Z-score's edge over PCA is small (0.01--0.02 in $r$) but statistically significant for all three metrics, with a larger, likewise significant gap to JointVAE (0.05--0.16); so the ordering in Table~\ref{tab:corr} is not an artifact of sampling noise, even though eyeballing the marginal intervals alone would suggest otherwise. Third, the observed ceiling near $r \approx 0.8$ versus WAR is consistent with equal weighting across dimensions and with the different construction of WAR; it should not be interpreted as a geometry-independent theoretical limit. Fourth, recomputing area under every axis permutation yields narrow correlation bands for all three encodings ($r$ vs.\ WAR: Z-score $[0.79, 0.82]$, PCA $[0.78, 0.81]$, JointVAE $[0.74, 0.76]$; comparable bands hold for wRC+ and OPS), although individual rankings can remain unstable. Notably, PCA is the most permutation-sensitive of the three encodings, both in absolute terms and as a fraction of its own correlation, so Z-score's advantage over the alternatives should be attributed to its higher, significant correlation with value metrics rather than to any special robustness to the pentagon's axis ordering.

\subsection{Qualitative Analysis: Profiles Behind Geometric Summaries}
Table~\ref{tab:top10} lists the top-10 season-rows by pentagon area (Z-score encoding) with their established metrics. Elite seasons by Ronald Acu\~na Jr., Shohei Ohtani, Aaron Judge, Juan Soto, Jose Ramirez, and Bobby Witt Jr.\ appear, but area also ranks profiles whose established offensive values differ substantially. Judge's 2024 vector is $(58, 79, 44, 37, 65)$ in (Contact, Power, Speed, Defense, Discipline) order---extreme Power with elite Discipline---while Soto's 2024 vector $(64, 66, 29, 44, 71)$ is carried by Discipline, and Witt's 2024 vector $(67, 60, 53, 60, 41)$ is more distributed. Ohtani's 2024 season $(59, 74, 72, 37, 48)$ pairs elite Power with elite Speed. These examples illustrate profile readability rather than external validation of the scores.

\begin{table}[t]
\centering
\caption{Top-10 season-rows by pentagon area (Z-score encoding) with established metrics.}
\label{tab:top10}
\scriptsize
\setlength{\tabcolsep}{3pt}
\begin{tabular}{lrrrr}
\toprule
Player (Season) & Area & WAR & wRC+ & OPS \\
\midrule
Ronald Acu\~na Jr.\ (2023) & 8644.6 &  9.2 & 171 & 1.012 \\
Shohei Ohtani (2024)       & 8089.1 &  9.0 & 180 & 1.036 \\
Aaron Judge (2022)         & 7794.7 & 11.1 & 206 & 1.111 \\
Juan Soto (2021)           & 7639.3 &  6.8 & 164 & 0.999 \\
Jose Ramirez (2021)        & 7568.0 &  6.4 & 140 & 0.893 \\
Aaron Judge (2024)         & 7523.9 & 11.3 & 220 & 1.159 \\
Bobby Witt Jr.\ (2024)     & 7349.6 & 10.5 & 169 & 0.977 \\
Nicky Lopez (2021)         & 7155.2 &  5.5 & 104 & 0.744 \\
Juan Soto (2024)           & 7128.0 &  8.3 & 181 & 0.989 \\
Bobby Witt Jr.\ (2025)     & 7102.5 &  8.0 & 130 & 0.852 \\
\bottomrule
\end{tabular}
\end{table}

The list also contains an instructive outlier: Nicky Lopez's 2021 season ranks eighth by area despite a merely league-average bat (104 wRC+, .744 OPS). His vector $(63, 30, 66, 71, 48)$ is the mirror image of a slugger's---elite Defense and Speed with well below-average Power---and the symmetric area rewards that balance. Notably, area here agrees with WAR (5.5, driven by defense and baserunning) far better than with the offense-only metrics, illustrating both what the equal-weight summary captures and what it overweights.

Figure~\ref{fig:radar} makes a profile contrast concrete. Bobby Witt Jr.\ (2024) is represented primarily through Contact, Power, and Speed, whereas Adley Rutschman's 2022 vector $(58, 45, 49, 63, 58)$ is stronger in Defense and Discipline. Their WAR values are not comparable (10.5 versus 5.6), so this pair is used only to illustrate divergent profile shapes and not as evidence that players with equal overall value are interchangeable.

\begin{figure}[t]
\centering
\includegraphics[width=\linewidth]{figures/fig2_radar_witt_vs_rutschman.png}
\caption{Radar-chart comparison of Bobby Witt Jr.\ (2024) and Adley Rutschman (2022). The players have different WAR values (10.5 and 5.6); the figure illustrates profile shape rather than matched overall value. Scores are 20--80-formatted relative values, not scout grades.}
\label{fig:radar}
\end{figure}

\subsection{Dimension-Delta Regression and Model Sensitivity}
\label{sec:roi-results}
Table~\ref{tab:cv} reports the $\Delta$Tool regression's $R^2$ under the three cross-validation schemes of Sec.~\ref{sec:roi-method}, for both the five-feature linear model and the ten-feature random forest.

\begin{table*}[t]
\centering
\caption{Cross-validation comparison for the dimension-delta regression (1{,}427 consecutive-season transitions). Pooled: shuffled 5-fold \texttt{KFold}, the original scheme. Player-grouped: \texttt{GroupKFold} keyed on player identity, so no player's transitions appear in both a fold's training and test split; additionally refitting the 20--80 scaling within each training fold changes every value by $<10^{-3}$ and is therefore not shown as a separate column. Forward-time holdout: trained on transitions ending $\leq$2023 ($n=707$), tested only on 2024$\to$2025 transitions ($n=362$).}
\label{tab:cv}
\small
\begin{tabular}{lcccccc}
\toprule
 & \multicolumn{2}{c}{Pooled KFold} & \multicolumn{2}{c}{Player-grouped} & \multicolumn{2}{c}{Forward-time holdout} \\
\cmidrule(lr){2-3} \cmidrule(lr){4-5} \cmidrule(lr){6-7}
Target & LR & RF & LR & RF & LR & RF \\
\midrule
$\Delta$wRC+ & 0.707 & 0.677 & 0.708 & 0.679 & 0.702 & 0.636 \\
$\Delta$WAR  & 0.704 & 0.673 & 0.703 & 0.677 & 0.692 & 0.644 \\
$\Delta$OPS  & 0.707 & 0.678 & 0.708 & 0.678 & 0.728 & 0.691 \\
\bottomrule
\end{tabular}
\end{table*}

Player-grouped cross-validation produces scores statistically indistinguishable from the pooled scheme across all three targets and both models (Table~\ref{tab:cv}), and refitting the 20--80 scaling on each training fold only changes nothing further, to the reported precision---consistent with linear-regression and random-forest predictions being invariant to the per-feature affine rescaling that step performs. This indicates the pooled estimate is not an artifact of within-player leakage for this model class and feature set. The forward-time holdout shows a modest but real gap for the random forest (0.68$\to$0.64 for both $\Delta$wRC+ and $\Delta$WAR), while the linear model remains comparatively stable and even improves for $\Delta$OPS; we read the forward-time result as the more honest estimate of forecasting a genuinely future season, and the random forest's gap as evidence it captures some era-specific structure that within-window, player-grouped cross-validation alone cannot detect. These values quantify association under each respective split, not causal explanation. In the fitted linear model, a one-point change in Power is associated with $+3.2$ wRC+, Contact with $+2.2$, and Discipline with $+0.7$, conditional on the other dimension changes. Speed and Defense coefficients are near zero for the offense-only wRC+ target and positive for WAR. Because the inputs and targets share underlying season statistics, these coefficients are descriptive partial associations rather than exchange rates.

Applying a synthetic $+5$-grade perturbation to each dimension of each 2025 player produces player-specific model responses. For Aaron Judge, the largest predicted response is assigned to Power ($+12.4$ wRC+); for Byron Buxton, it is assigned to Contact ($+14.0$ wRC+). These examples show how the model can be inspected, but there is no observed ground truth establishing that the selected dimension is the player's optimal development target. The perturbations may also leave the joint data manifold. We therefore report them as sensitivity examples, not counterfactual outcomes, ROI, or prescriptions.

\begin{figure}[t]
\centering
\includegraphics[width=\linewidth]{figures/roi_coefficients.png}
\caption{Linear-regression coefficients associating concurrent year-over-year dimension changes with changes in value metrics. Coefficients are descriptive associations, not causal exchange rates.}
\label{fig:roi}
\end{figure}

\subsection{Sensitivity Analysis}
Because pentagon area sums adjacent products (Eq.~\ref{eq:area}), it is not permutation-invariant. Population-level correlations vary only within $\pm 0.02$ across orderings, but individual rankings can change substantially as within-player dimension variance increases. Thus, the ``spiky'' profiles for which the full vector is most informative are also those for which radar-chart area is least stable. This finding limits area as a comparison metric and supports using it only as an auxiliary geometric descriptor.

\subsection{Exploratory External Comparison}
\label{sec:external-pilot}
As a small, qualitative check against external construct validity, we compared six recently graduated top prospects' first full-plate-appearance-season Baseball2Vector vectors against publicly available, individually cited pre-MLB scouting grades (Baseball America; exact source URLs are released with the repository). Table~\ref{tab:pilot} restricts the comparison to the three dimensions with a plausible cross-taxonomy mapping (Contact$\leftrightarrow$Hit, Power$\leftrightarrow$Power, Speed$\leftrightarrow$Run); Defense and Plate Discipline are excluded from the table because public sources either split defense into separate Field and Arm grades that Baseball2Vector does not, or, for Discipline, do not report a 20--80 grade at all. The two scales are not numerically interchangeable: public grades are calibrated against the full population of minor-league prospects, whereas Baseball2Vector's scale is calibrated only against players who already reached MLB regular status, so we discuss relative profile shape rather than literal number matching.

Speed showed the most consistent agreement---for example, Corbin Carroll's Speed score of 78.8 closely tracked a scouted 80-grade run tool, the closest match in the panel. Power showed the most consistent disagreement, reading systematically lower in Baseball2Vector than the scouted raw-power grade (e.g.\ Julio Rodr\'iguez, Elly De La Cruz), plausibly because scouted power grades often describe raw batting-practice ceiling while Baseball2Vector's Power dimension is built entirely from in-game production statistics---a genuine difference in what is measured, not necessarily an error. With six players and two non-commensurable scales, this pilot is hypothesis-generating only and does not establish external validity (Limitations, below).

\begin{table}[t]
\centering
\caption{Exploratory pilot: Baseball2Vector (Z-score) vs.\ public pre-MLB scouting grades (Baseball America), restricted to the three dimensions with a plausible cross-taxonomy mapping. B2V values use each player's first full-PA MLB season; public grades predate MLB debut. Not a validation sample ($n=6$); see Sec.~\ref{sec:external-pilot}.}
\label{tab:pilot}
\scriptsize
\setlength{\tabcolsep}{3pt}
\begin{tabular}{lccc}
\toprule
Player (B2V season) & Contact / Hit & Power / Power & Speed / Run \\
\midrule
Witt Jr.\ (2022)      & 52.9 / 60 & 46.3 / 60 & 64.0 / 60 \\
Rodr\'iguez (2022)    & 49.9 / 70 & 58.6 / 80 & 54.5 / --- \\
Rutschman (2022)      & 58.1 / 70 & 44.5 / --- & 48.5 / --- \\
Henderson (2023)      & 50.6 / 60 & 54.4 / 70 & 55.5 / 60 \\
Carroll (2023)        & 58.5 / 60 & 49.2 / 60 & 78.8 / 80 \\
De La Cruz (2023)     & 44.1 / 42 & 51.7 / 70 & 60.7 / 70 \\
\bottomrule
\end{tabular}
\end{table}

\section{Conclusion and Limitations}
\label{sec:conclusion}

We presented Baseball2Vector, a mapping from public season statistics to an interpretable five-dimensional profile reported on a 20--80-formatted relative scale. Across 2{,}309 hitter-seasons (2021--2025), Z-score averaging showed the strongest descriptive alignment with established value metrics ($r=0.81$ vs.\ WAR). Qualitative cases illustrate how the vector exposes dimension-level differences that its geometric area discards. Dimension-delta regression produced a pooled five-fold $R^2$ of 0.71 for $\Delta$wRC+, a figure that player-grouped cross-validation shows is not primarily a within-player leakage artifact (Sec.~\ref{sec:roi-results}), though it remains an observational, shared-input association rather than evidence of causal development effects.

\paragraph{Limitations.}
The central limitation is circularity: several input statistics also contribute directly or indirectly to WAR, wRC+, or OPS, so high alignment partly reflects re-expression of shared information rather than independent validation. The proposed dimensions also differ from the canonical Hit--Power--Run--Field--Throw framework; a small pilot comparison against publicly available scouting grades for six recently graduated prospects (Sec.~\ref{sec:external-pilot}) is suggestive but far too small to establish construct validity. Min--max mapping produces cohort-dependent 20--80-formatted scores rather than classical scouting grades. Within and across dimensions, correlated inputs can double-count related signals; defensive composites may also mix fielding performance, position, playing time, and opportunity. Pentagon area is equal-weighted and ordering-sensitive. Repeated seasons and overlapping transitions from the same player create dependence; player-grouped and forward-time cross-validation (Sec.~\ref{sec:roi-results}) indicate the pooled dimension-delta estimate is not primarily a leakage artifact for the model classes considered, though the random forest's forward-time degradation suggests some era-specific structure that within-window grouped cross-validation alone cannot detect. Finally, the synthetic $+5$ perturbations are correlational model-sensitivity probes that inherit survivorship, role-change, age, injury, and playing-time confounds. Coverage is restricted to MLB hitters with $\geq$100 PA and excludes pitchers, minor leaguers, and most short-sample players. Exact reproduction additionally requires a versioned public repository containing the extraction date, processed data, preprocessing code, package versions, random seeds, and fold assignments.

\paragraph{Future directions.}
The next priority is external validation at scale rather than a more complex encoder. First, the pilot comparison of Sec.~\ref{sec:external-pilot} should be extended well beyond six players, ideally with a scouting organization's cooperation, before any construct-validity claim stronger than a qualitative pattern is made. Second, while transition models have now been re-evaluated with player-grouped cross-validation, fold-refit preprocessing, and a forward-time holdout (Sec.~\ref{sec:roi-results}), raw-statistic, previous-season, and no-change baselines with uncertainty intervals remain to be reported, and the random forest's forward-time degradation warrants direct investigation into whether era-specific offensive-environment shifts are the cause. Third, synthetic perturbations should be replaced or complemented by evaluation against subsequently realized player changes, with age, role, position, playing time, and injury status modeled explicitly. Only after these tests should the representation be considered for predictive or development-support applications.

%%%%%%%%% REFERENCES
\begin{thebibliography}{99}

\bibitem{mcdaniel2014}
K. McDaniel.
\newblock Scouting explained: The 20--80 scouting scale.
\newblock \emph{FanGraphs Baseball}, 2014. \url{https://blogs.fangraphs.com/scouting-explained-the-20-80-scouting-scale/}.

\bibitem{alcorn2016}
M. A. Alcorn.
\newblock 2vec: Statistic-free talent modeling with neural player embeddings.
\newblock In \emph{MIT Sloan Sports Analytics Conference}, 2016.

\bibitem{heaton2021}
C. Heaton and P. Mitra.
\newblock Learning to describe player form in the MLB.
\newblock \emph{arXiv preprint arXiv:2109.05280}, 2021.

\bibitem{fangraphsSampleSize}
FanGraphs Sabermetrics Library.
\newblock Sample size.
\newblock \url{https://library.fangraphs.com/principles/sample-size/}.

\bibitem{mikolov2013}
T. Mikolov, K. Chen, G. Corrado, and J. Dean.
\newblock Efficient estimation of word representations in vector space.
\newblock {\em arXiv preprint arXiv:1301.3781}, 2013.

\bibitem{kingma2014}
D. P. Kingma and M. Welling.
\newblock Auto-encoding variational Bayes.
\newblock In {\em ICLR}, 2014.

\end{thebibliography}

\end{document}
